\subsection{\texorpdfstring{\(p\)-radical extensions}{p-radical extensions}}

DEFINITION 2. --- Let E be an extension of a field K. We shall say that E is p-radical \(^{1}\) (over K) if every element of E is p-radical over K. When this is so, E is said to be of finite height when the set of heights of elements of E is bounded above and we call the height of E the maximum of the heights of its elements.

We note that every \(p\) -radical extension of a perfect field (in particular of a field of characteristic 0) is trivial.

Let K be a field.. The p-radical extensions of height 0 of K are the trivial extensions. Every p-radical extension of K is algebraic. If E is a p-radical extension of K and F a p-radical extension of E, then F is a p-radical extension of K: for, given any \(x \in F\) , there exists an integer \(m \geqslant 0\) such that \(x^{p^{m}} \in E\) and an integer \(n \geqslant 0\) such that \((x^{p^{m}})^{p^{n}} \in \mathbf{K}\) , that is, \(x^{p^{m+n}} \in K\) .

PROPOSITION 2. --- Let E be an extension of a field K. For every integer \(n \geqslant 0\) let \(E_{n}\) be the set of elements of E which are p-radical of height \(\leqslant n\) over K, and let \(E_{\infty}\) be the set of all elements of E that are p-radical over K. Then \((\mathrm{E}_{n})_{n \geqslant 0}\) is an ascending sequence of subextensions of E whose union is \(E_{\infty}\) , and \(E_{\infty}\) is the largest p-radical extension of K contained in E.

For each integer \(n \geqslant 0\) the set \(E_{n}\) consists of all elements x of E such that \(x^{p^{n}} \in K\) ; since the mapping \(x \mapsto x^{p^{n}}\) is an endomorphism of the field E, we deduce that \(E_{n}\) is a subextension of E. The sequence \((\mathrm{E}_{n})_{n \geqslant 0}\) is ascending, with union \(E_{\infty}\) and hence \(E_{\infty}\) is a subextension of E (V, p.~11, Prop. 3). It is clear that \(E_{\infty}\) is a p-radical extension of K and that \(E_{\infty}\) contains every subextension of E which is p-radical over K.

COROLLARY. --- If an extension E of a field K is generated by a set of p-radical elements over K, then it is p-radical over K.

For \(E_{\infty}\) is a subextension of E and by hypothesis \(\mathrm{E} = \mathrm{K}(\mathrm{E}_{\infty})\) whence \(E = E_{\infty}\) ; so E is a p-radical extension of K.

Under the conditions of Proposition 2, \(E_{\infty}\) is said to be the relative p-radical closure of K in E.

We shall apply Prop. 2 particularly to the case where E is an algebraically closed extension of K; then E is perfect and we have \(E_{n} = K^{p^{-n}}\) for all \(n \geqslant 0\) . In this case we denote by \(K^{p^{-\infty}}\) the set of elements of E that are p-radical over K; it is the subfield of E which is the union of the ascending sequence \((\mathbf{K}^{p^{-n}})_{n \geqslant 0}\) of subfields of E. By Prop. 2, \(K^{p^{-\infty}}\) is the largest subextension of E which is p-radical over K; by Prop. 3 of V, p.~6, \(K^{p^{-\infty}}\) is a perfect closure of K and it is also the smallest perfect subfield of E containing K. When K is perfect we clearly have \(K = K^{p^{-n}} = K^{p^{-\infty}}\) for all n.~If K is imperfect, we have \(K \neq K^{p}\) , whence

\[
\mathbf {K} ^ {p ^ {- n}} \neq (\mathbf {K} ^ {p}) ^ {p ^ {- n}} = \mathbf {K} ^ {p ^ {- (n - 1)}}
\]

for \(n \geqslant 1\) ; the subfields \(\mathbf{K}^{p^{-n}}\) of \(E\) are then pairwise distinct and \(\mathbf{K}^{p^{-\infty}}\) is an algebraic extension of infinite degree of \(K\) .

PROPOSITION 3. --- Let K be a field, E an extension field of K which is p-radical over K and u a homomorphism of K into a perfect field F. Then there exists a unique homomorphism v of E into F extending u.

Let \(E_n\) be the set of elements of \(E\) which are \(p\) -radical of height \(\leqslant n\) over \(K\) . By Prop. 2 the field \(E\) is the union of the ascending sequence \((E_n)_{n \geqslant 0}\) of subfields. Let \(v\) be a homomorphism of \(E\) into \(F\) extending \(u\) ; for any \(x \in E_n\) we have \(x^{p^n} \in K\) , whence \(v(x)^{p^n} = v(x^{p^n}) = u(x^{p^n})\) ; we thus have \(v(x) = u(x^{p^n})^{p^{-n}}\) for all \(n \geqslant 0\) and all \(x \in E_n\) , whence the uniqueness of the extension of \(u\) to \(E\) .

Let \(n\) be a positive integer; for any \(x \in \mathrm{E}_n\) we have \(x^{p^n} \in \mathbf{K}\) and since \(F\) is perfect, we can define an element \(v_n(x)\) of \(F\) by \(v_n(x) = u(x^{p^n})^{p^{-n}}\) . It is clear that \(v_n\) is a homomorphism of \(E_n\) into \(F\) , that \(v_0 = u\) and that \(v_{n+1}\) induces \(v_n\) on \(E_n\) . Hence there exists a homomorphism \(v\) of \(E\) into \(F\) inducing \(v_n\) on \(E_n\) for all \(n \geqslant 0\) and, in particular, coinciding with \(v_0 = u\) on \(E_0 = K\) .

COROLLARY. --- For an extension E of a field K to be a perfect closure of K it is necessary and sufficient that it should be a p-radical extension of K and that the field E should be perfect.

The Corollary is trivial when p = 1; suppose then that \(p \neq 1\) . The stated conditions are necessary by V, p.~5, Th. 3; they are sufficient by Prop. 3.

PROPOSITION 4. --- Let E be a p-radical extension of finite degree of a field K. Then {[}E : K{]} is a power of the characteristic exponent p of K.

Since E is a p-radical extension of finite degree of K, there are elements \(a_{1}, \ldots, a_{m}\) of E, p-radical over K, such that \(\mathrm{E} = \mathrm{K}(a_{1}, \ldots, a_{m})\) . Let i be in the range 1 to m; since \(a_{i}\) is a fortiori p-radical over \(\mathrm{K}(a_{1}, \ldots, a_{i-1})\) , the degree

\[
n _ {i} = \left[ \mathrm{K} (a _ {1}, \dots , a _ {i}): \mathrm{K} (a _ {1}, \dots , a _ {i - 1}) \right]
\]

is a power of \(p\) (V, p.~24, Prop. 1). We have \([\mathrm{E}:\mathbf{K}] = n_1\ldots n_m\) by V, p.~18, Th. 2, whence the result.

